% !TeX program = xelatex
% 从 slide/ 目录编译；逐页讲述时间与代码出处见 README.md。
\documentclass[aspectratio=169,11pt]{beamer}
\usepackage[UTF8,fontset=fandol]{ctex}
\usepackage{graphicx,tikz,booktabs,tabularx,listings}
\usetikzlibrary{positioning,arrows.meta}
\definecolor{accent}{HTML}{245777}
\definecolor{ink}{HTML}{263440}
\definecolor{muted}{HTML}{657582}
\definecolor{pale}{HTML}{F0F3F5}
\setbeamercolor{normal text}{fg=ink,bg=white}
\setbeamercolor{structure}{fg=accent}
\setbeamercolor{frametitle}{fg=accent}
\setbeamercolor{block title}{fg=accent,bg=pale}
\setbeamercolor{block body}{fg=ink,bg=pale}
\setbeamerfont{frametitle}{size=\Large,series=\bfseries}
\setbeamertemplate{navigation symbols}{}
\setbeamertemplate{itemize item}{\small\textbullet}
\setbeamertemplate{footline}{\hspace{0.55cm}\textcolor{muted}{\scriptsize AI Photo Coach\quad 端云协同智能摄影辅助}\hfill\textcolor{muted}{\scriptsize\insertframenumber/\inserttotalframenumber}\hspace{0.55cm}\vspace{0.22cm}}
\addtobeamertemplate{frametitle}{}{\vspace{0.12cm}}
\setbeamersize{text margin left=0.65cm,text margin right=0.65cm}
\setlength{\parskip}{0.3em}
% 所有素材均位于 assets/；缺失时保留尺寸并显示文件名。
\newcommand{\asset}[3]{\IfFileExists{assets/#1}{\begin{minipage}[c][#3][c]{#2}\centering\includegraphics[width=#2,height=#3,keepaspectratio]{assets/#1}\end{minipage}}{\begingroup\setlength{\fboxsep}{0pt}\colorbox{pale}{\begin{minipage}[c][#3][c]{#2}\centering\color{muted}\tiny 素材待补充\par\ttfamily\expandafter\detokenize\expandafter{#1}\end{minipage}}\endgroup}}
\newcommand{\qwenfile}{logo_qwen.png}
\IfFileExists{assets/logo_qwen.png}{}{\IfFileExists{assets/log_qwen.png}{\renewcommand{\qwenfile}{log_qwen.png}}{}}
\newcommand{\takeaway}[1]{\vspace{0.2cm}\begin{beamercolorbox}[sep=0.13cm,wd=\linewidth]{block body}\centering\small\color{accent}#1\end{beamercolorbox}}
\tikzset{box/.style={draw=accent!30,fill=pale,rounded corners=2pt,align=center,minimum height=2.65cm,text width=2.8cm,font=\footnotesize},arr/.style={-{Stealth[length=2mm]},thick,accent}}
\lstset{basicstyle=\ttfamily\fontsize{9.5}{10.5}\selectfont,keywordstyle=\color{accent}\bfseries,commentstyle=\color{muted},stringstyle=\color{accent},backgroundcolor=\color{pale},frame=none,breaklines=true,columns=fullflexible,keepspaces=true,showstringspaces=false,xleftmargin=0.15cm,xrightmargin=0.15cm,aboveskip=0.12cm,belowskip=0.12cm}
\begin{document}

\begin{frame}[plain]
\vspace{0.65cm}
\begin{center}
{\color{accent}\fontsize{25}{32}\selectfont\bfseries AI Photo Coach\par 智能摄影辅助工具\par}
{\color{muted}\normalsize 基于端云协同的智能人像摄影辅助系统\par}
\vspace{0.25cm}
{\large 2026短\quad 图灵实践课II 课程考核项目答辩\par}
\end{center}
\vspace{0.4cm}
\vfill
\begin{columns}[b,onlytextwidth]
\column{0.48\textwidth}
\small 小组成员：郭名扬\quad 刘泽宇\par
\vspace{0.2cm}\scriptsize\color{muted}2026 年 9 月 10 日\par
\column{0.49\textwidth}\raggedleft
\asset{logo_orangepi.png}{2.3cm}{1.3cm}\hspace{0.3cm}
\asset{logo_ascend.png}{2.3cm}{1.3cm}\hspace{0.3cm}
\asset{logo_yolo.png}{2.3cm}{1.3cm}\hspace{0.3cm}
\asset{\qwenfile}{2.3cm}{1.3cm}
\end{columns}\vspace{0.5cm}
\end{frame}

\begin{frame}{项目背景}
\begin{columns}[T,onlytextwidth]
\column{0.47\textwidth}
\asset{logo_yolo.png}{2.3cm}{1.3cm}
\begin{block}{拍摄中：实时指导}
\begin{itemize}
\item 实时、低延迟的构图反馈
\item 明确告诉用户如何调整取景
\item 关注位置、留白、主体大小
\end{itemize}
\end{block}
\column{0.47\textwidth}
\asset{\qwenfile}{2.3cm}{1.3cm}
\begin{block}{拍摄后：进一步分析}
\begin{itemize}
\item 理解光线与背景干扰
\item 分析姿态与画面表达
\item 给出有原因的摄影建议
\end{itemize}
\end{block}
\end{columns}
\vspace{0.3cm}
\small YOLO 负责检测人物，OpenCV 可以计算位置和亮度，但难以判断光线质量、背景干扰和人物状态。VLM 可以找出更深层的语义，但会带来更大的延迟。
\takeaway{因此可以把实时分析放在香橙派，把拍照后的深入分析交给云端。}
\end{frame}

\begin{frame}{系统工作流程}{拍摄时：香橙派快速分析}
\centering
\begin{tikzpicture}[node distance=0.32cm]
\node[box] (a) {\asset{Vue.png}{1.7cm}{1cm}\\Vue 3\\Canvas\\Camera API\\Browser};
\node[box,right=of a] (b) {\asset{logo_orangepi.png}{2.6cm}{1.7cm}\\Orange Pi\\AI Pro 20T\\Nginx Web 入口};
\node[box,right=of b] (c) {\asset{logo_fastapi.png}{2.6cm}{1.7cm}\\FastAPI\\API 请求处理\\边缘分析服务};
\node[box,right=of c] (d) {\asset{logo_ascend.png}{1.7cm}{1cm}\\Ascend 310B1\\YOLOv5s OM\\ACLLite\\OpenCV\\构图规则};
\draw[arr] (a)--(b);\draw[arr] (b)--(c);\draw[arr] (c)--(d);
\end{tikzpicture}
\vspace{0.45cm}
\begin{itemize}
\item 浏览器把摄像头画面压缩成最长边 640px 的 JPEG。
\item 香橙派完成人物检测，并计算位置、留白、人物大小和亮度。
\item 网页拿到结果后，更新人物框、构图参考和拍摄建议。
\end{itemize}
\takeaway{这部分重点满足低延迟的需求，让用户可以在拍摄过程中快速调整。}
\end{frame}

\begin{frame}{系统工作流程}{拍照后：云端按需深入分析}
\centering
\begin{tikzpicture}[node distance=0.32cm]
\node[box] (a) {\asset{logo_fastapi.png}{1.7cm}{1cm}\\FastAPI\\照片 + 构图数据};
\node[box,right=of a] (b) {\asset{ssh.png}{1.7cm}{1cm}\\SSH Tunnel\\本地端口转发};
\node[box,right=of b] (c) {\asset{logo_modelarts.png}{1.7cm}{1cm}\\华为云 ModelArts\\Ascend 910B3};
\node[box,right=of c] (d) {\asset{\qwenfile}{1.7cm}{1cm}\\Qwen3-VL\\8B-Instruct};
\draw[arr] (a)--(b);\draw[arr] (b)--(c);\draw[arr] (c)--(d);
\end{tikzpicture}
\vspace{0.4cm}
\begin{itemize}
\item 用户点击“AI 深度分析”后，系统才会上传照片和 YOLO / CV 数据。
\item 香橙派通过 SSH 隧道调用 ModelArts 上的 Qwen3-VL 服务。
\item Cloudflare Quick Tunnel 为手机演示提供临时 HTTPS 地址。
\end{itemize}
\takeaway{这部分是让 VLM 根据理解到的深层语义返回更详细的建议。}
\end{frame}

\begin{frame}{拍摄时：实时给出调整建议}
\begin{columns}[c,onlytextwidth]
\column{0.53\textwidth}\centering
\asset{screenshot_live_advice.png}{\linewidth}{7cm}
\column{0.43\textwidth}
\centering
\begin{tikzpicture}[node distance=0.13cm, every node/.style={font=\small}]
\node (a) {Camera frame};
\node[below=of a] (b) {Canvas：640px JPEG};
\node[below=of b] (c) {\texttt{/api/analyze-fast}};
\node[below=of c] (d) {YOLO bbox};
\node[below=of d] (e) {构图规则};
\node[below=of e] (f) {实时建议};
\foreach \s/\t in {a/b,b/c,c/d,d/e,e/f}{\draw[arr] (\s)--(\t);}
\end{tikzpicture}
\vspace{0.2cm}
\begin{block}{反馈内容}
\small 人物位置\quad 上方留白\par 主体占比\quad 亮度统计
\end{block}
\end{columns}
\end{frame}

\begin{frame}{构图建议如何生成}
\begin{columns}[c,onlytextwidth]
\column{0.54\textwidth}\centering
\asset{screenshot_photo_analysis.png}{\linewidth}{6.35cm}
\column{0.42\textwidth}
{\color{accent}\normalsize bbox $\rightarrow$ 几何指标}
\footnotesize
\begin{align*}
\text{center}_x &= \frac{x_1+x_2}{2W}\\
\text{headroom} &= \frac{y_1}{H}\\
\text{area ratio} &= \frac{(x_2-x_1)(y_2-y_1)}{WH}
\end{align*}
\vspace{-0.4cm}
\begin{itemize}
\item 参考三分线，也允许大面积人像居中。
\item 位置、上方留白、人物大小组成参考分。
\item 偏差再变成“向左”“拉远”等具体建议。
\end{itemize}
\takeaway{这个分数代表几何构图，不代表照片整体的美观性；\par 同时亮度单独显示，不计入构图分。}
\end{columns}
\end{frame}

\begin{frame}{拍照后：Qwen3-VL 深入分析}
\begin{columns}[c,onlytextwidth]
\column{0.53\textwidth}\centering
\asset{screenshot_deep_analysis.png}{\linewidth}{7cm}
\column{0.43\textwidth}
\begin{block}{输入}
\small 当前照片 + YOLO / CV 数据
\end{block}
\centering $\downarrow$\par
{\color{accent}\large Qwen3-VL-8B-Instruct}\par
$\downarrow$
\begin{block}{返回结果}
\small scene / lighting / background\par
pose / visual\_expression / suggestions
\end{block}
\raggedright\small 在已有构图数据的基础上，再分析场景、光线、背景和姿态，给出更具体的调整建议。
\end{columns}
\end{frame}

\begin{frame}[fragile]{关键代码}{A · YOLO 人物框如何回到原图}
\setlength{\parskip}{0pt}
\footnotesize YOLO 输出的是 640×640 letterbox 坐标；下面将人物框还原到原图，并转换为网页可用的比例。
\begin{lstlisting}[language=Python]
x1 = (model_x1 - pad_x) / scale
y1 = (model_y1 - pad_y) / scale
x2 = (model_x2 - pad_x) / scale
y2 = (model_y2 - pad_y) / scale
# Clip all coordinates to the original image boundary
x1 = int(max(0, min(x1, original_width - 1)))
bbox_norm = {
    "x1": round(x1 / original_width, 4),
    "y1": round(y1 / original_height, 4),
    # x2 / y2 are handled in the same way
}
\end{lstlisting}
\takeaway{先减去补边，再除以缩放比例；\par 归一化坐标让人物框在不同屏幕尺寸下仍能对准照片。}
\scriptsize\color{muted}来源：\texttt{edge-backend/detector.py}，\texttt{AscendYoloDetector.detect()}
\end{frame}

\begin{frame}[fragile]{关键代码}{B · 如何把人物框变成构图建议}
\setlength{\parskip}{0pt}
\small 项目表达式节选，合并换行；\texttt{cx\_ratio} 即归一化的 center\_x。
\begin{lstlisting}[language=Python]
cx = (x1 + x2) / 2
cx_ratio = cx / W
person_width = max(0, x2 - x1)
person_height = max(0, y2 - y1)
person_area = person_width * person_height
image_area = W * H
area_ratio = person_area / image_area
headroom_ratio = y1 / H
# ...
score = round(position_score * 0.45
              + headroom_score * 0.30
              + subject_size_score * 0.25)
\end{lstlisting}
\takeaway{坐标先转成比例，换一台手机仍能使用同一套规则；\par 每条建议也能找到对应的计算依据。}
\scriptsize\color{muted}来源：\texttt{edge-backend/composition.py}，\texttt{analyze\_composition()}
\end{frame}

\begin{frame}[fragile]{关键代码}{C · 如何把照片和构图数据交给 VLM}
\setlength{\parskip}{0pt}
\small 边缘端请求节选：图像和指标通过同一请求送达云端。
\begin{lstlisting}[language=Python]
response = self._client.post(
    self.api_url,
    files={"image": ("frame.jpg", image_bytes, "image/jpeg")},
    data={"cv_metrics": json.dumps(metrics, ensure_ascii=False)},
)
\end{lstlisting}
\small 云端消息节选：\texttt{build\_prompt(metrics)} 写入指标与 JSON 输出约束。
\begin{lstlisting}[language=Python]
"content": [
    {"type": "image", "image": image},
    {"type": "text", "text": build_prompt(metrics)},
],
\end{lstlisting}
\takeaway{提示词约定返回 JSON；服务端负责解析和整理字段，\par 出错时保留原有的 YOLO / CV 结果。}
\scriptsize\color{muted}来源：\texttt{edge-backend/services/vlm\_client.py}；\texttt{cloud-vlm/app.py}
\end{frame}

\begin{frame}{实际运行效果}
\vspace{1.1cm}
\centering\Large\color{accent}
实时指导\quad $\rightarrow$\quad 拍摄\quad $\rightarrow$\quad 照片诊断\quad $\rightarrow$\quad AI 深度分析
\vfill
% 视频外部播放，预留路径 assets/demo.mp4；本页不嵌入媒体。
\end{frame}

\begin{frame}{遇到的问题与解决办法}
\renewcommand{\arraystretch}{1.65}
\small
\begin{tabularx}{\linewidth}{@{}p{3.5cm}X@{}}
\toprule
\textbf{问题} & \textbf{解决方法}\\
\midrule
大图上传慢 & 浏览器端压缩为最长边 640px 的 JPEG\\
实时请求堆积 & 串行请求 + 互斥状态，避免重复启动\\
只靠构图规则，建议偏机械 & 明确标为“构图参考”，再由 VLM 补充画面理解\\
云端慢或失败 & 快速分析继续可用，深度分析降级并提示重试\\
手机摄像头 HTTPS 要求 & Cloudflare Quick Tunnel 提供演示入口\\
\bottomrule
\end{tabularx}
\vspace{0.3cm}
\takeaway{实时指导和深度分析互不绑死：云端暂时不可用，基本拍摄功能还能继续。}
\end{frame}

\begin{frame}{项目总结与下一步}
\begin{columns}[T,onlytextwidth]
\column{0.31\textwidth}
\begin{block}{已完成}
\small
实时指导、照片诊断\par
边缘 YOLO\par
云端 Qwen3-VL\par
端云连接、演示网页
\end{block}
\column{0.31\textwidth}
\begin{block}{当前限制}
\small
面向单人人像\par
构图分并非通用审美分\par
VLM 有生成不确定性
\end{block}
\column{0.31\textwidth}
\begin{block}{下一步}
\small
固定公网入口\par
%异步深度任务\par
更多构图规则\par
真实数据评测集
\end{block}
\end{columns}
\vspace{0.65cm}
\centering\large\color{accent}香橙派侧负责实时分析反馈，云端模型负责回馈更详细的建议。
\end{frame}

\begin{frame}
\vfill
\begin{center}
    {\Huge \textbf{Thank You}}

    \vspace{1.8cm}

    {\small 小组成员：郭名扬\quad 刘泽宇\par}
\end{center}
\end{frame}
\end{document}
